\documentclass[a4paper]{ctexart}

\usepackage{amsmath}
\usepackage{amssymb}
\usepackage{amsthm}
\usepackage[top=2.4cm,bottom=2.6cm,left=2.5cm,right=2.5cm]{geometry}
\usepackage[colorlinks=true,linkcolor=blue,pdfpagelabels=false]{hyperref}
\usepackage{bookmark}
\usepackage{enumitem}
\usepackage{mathtools}

\linespread{1.2}

\title{Algebra III 对应 Atiyah 习题详解}
\author{}
\date{2026-06-29}

\begin{document}
\maketitle
\tableofcontents

\section*{第0章：说明}
\phantomsection
\addcontentsline{toc}{section}{第0章：说明}

本文把 \textit{Algebra III} 中明确标注为
\[
\texttt{Exercise x.x.x. [2] Chapter \(\cdots\), exercises \(\cdots\)}
\]
的题目，按 Atiyah--Macdonald 的原章节统一摘出，并给出比速记版更详细的证明。

文中若直接调用 \textit{Algebra III} 的定理或命题，会尽量标明原书位置，例如
\[
\text{(III.Prop 1.2.13)},\quad \text{(III.Cor 2.2.6)},\quad \text{(III.Thm 8.2.2)}.
\]
这里只标出真正参与论证的结果；完全初等的张量、局部化或有限生成模计算不再强行附会引用。

对应关系如下：
\begin{itemize}
  \item Algebra III 1.2.18 对应 Atiyah 第2章习题 5, 8, 10, 12, 13。
  \item Algebra III 2.2.12 对应 Atiyah 第3章习题 2, 3, 15, 19, 21, 22。
  \item Algebra III 6.2.13 对应 Atiyah 第4章习题 2, 4, 5。
  \item Algebra III 4.3.7 对应 Atiyah 第5章习题 1, 4, 5, 7, 12, 13, 14, 15。
  \item Algebra III 5.3.17 对应 Atiyah 第6章习题 1, 2；第7章习题 1, 2, 12；第8章习题 2, 3, 4。
  \item Algebra III 7.2.12 对应 Atiyah 第9章习题 2, 4, 7, 9。
  \item Algebra III 8.5.15 对应 Atiyah 第10章习题 3, 4, 6, 7, 9。
\end{itemize}

\section{第2章：Modules}

\begin{enumerate}[leftmargin=*, label=\textbf{2.\arabic*.}]
\item 证明 \(A[x]\) 作为 \(A\)-代数是 flat。

\textbf{证明。}
把 \(A[x]\) 只看成 \(A\)-模，则
\[
A[x]=\bigoplus_{n\ge 0}Ax^n.
\]
右端是以 \(\{x^n\}_{n\ge 0}\) 为基的自由 \(A\)-模，因此 \(A[x]\) 是自由模。自由模一定是 flat 模，所以 \(A[x]\) flat。

\setcounter{enumi}{7}
\item 证明：
\begin{enumerate}[label=(\roman*)]
\item 若 \(M,N\) 都是 flat 的 \(A\)-模，则 \(M\otimes_A N\) 也是 flat。
\item 若 \(B\) 是 flat 的 \(A\)-代数，且 \(N\) 是 flat 的 \(B\)-模，则 \(N\) 也是 flat 的 \(A\)-模。
\end{enumerate}

\textbf{证明。}
\begin{enumerate}[label=(\roman*)]
\item 按定义，只需证函子 \(-\otimes_A(M\otimes_A N)\) 正合。由张量积结合律，
\[
X\otimes_A(M\otimes_A N)\cong (X\otimes_A M)\otimes_A N
\]
对每个 \(A\)-模 \(X\) 自然成立。因此
\[
-\otimes_A(M\otimes_A N)\cong \bigl((-\otimes_A M)\otimes_A N\bigr).
\]
由于 \(M\) flat，函子 \(-\otimes_A M\) 正合；由于 \(N\) flat，函子 \(-\otimes_A N\) 也正合。两个正合函子的复合仍正合，故 \(M\otimes_A N\) flat。

\item 对任意 \(A\)-模 \(X\)，有自然同构
\[
X\otimes_A N \cong (X\otimes_A B)\otimes_B N.
\]
这里先把 \(X\) 经标量扩张变成 \(B\)-模，再与 \(N\) 张量。

现在若
\[
0\to X'\to X\to X''\to 0
\]
是 \(A\)-模短正合列，则由于 \(B\) 对 \(A\) flat，得到
\[
0\to X'\otimes_A B\to X\otimes_A B\to X''\otimes_A B\to 0
\]
仍正合。再由于 \(N\) 对 \(B\) flat，对上式再张量 \(N\) 得
\[
0\to (X'\otimes_A B)\otimes_B N
\to (X\otimes_A B)\otimes_B N
\to (X''\otimes_A B)\otimes_B N
\to 0.
\]
借助前述自然同构，这正是
\[
0\to X'\otimes_A N\to X\otimes_A N\to X''\otimes_A N\to 0.
\]
故 \(-\otimes_A N\) 正合，即 \(N\) 对 \(A\) flat。
\end{enumerate}

\setcounter{enumi}{9}
\item 设 \(\mathfrak a\subseteq J(A)\)，\(u:M\to N\) 是 \(A\)-模同态，且 \(N\) 有限生成。若诱导映射
\[
M/\mathfrak aM\to N/\mathfrak aN
\]
满射，证明 \(u\) 满射。

\textbf{证明。}
设 \(C=\operatorname{coker}(u)\)。因为 \(N\) 有限生成，所以其商模 \(C\) 也是有限生成。

由右正合性，
\[
(N/u(M))/\mathfrak a(N/u(M)) \cong C/\mathfrak a C
\]
而题设说 \(M/\mathfrak aM\to N/\mathfrak aN\) 满射，这恰好等价于
\[
C/\mathfrak a C=0.
\]
于是 \(\mathfrak a C=C\)。又 \(\mathfrak a\subseteq J(A)\)，由 Nakayama 引理，可直接应用 \textit{Algebra III} 的 (III.Prop 1.2.13)，推出 \(C=0\)。因此 \(\operatorname{coker}(u)=0\)，即 \(u\) 满射。

\setcounter{enumi}{11}
\item 设 \(M\) 有限生成，\(\varphi:M\to A^n\) 满射。证明 \(\ker\varphi\) 有限生成。

\textbf{证明。}
令 \(e_1,\dots,e_n\) 为 \(A^n\) 的标准基。由 \(\varphi\) 满射，对每个 \(i\) 可取 \(u_i\in M\) 使得
\[
\varphi(u_i)=e_i.
\]
记
\[
F=\sum_{i=1}^n Au_i\subseteq M.
\]
则 \(\varphi|_F:F\to A^n\) 满射；而 \(F\) 由 \(n\) 个元生成，并且 \(u_i\) 在 \(A^n\) 中的像正好是标准基，故 \(\varphi|_F\) 实际上是同构。于是 \(M=F+\ker\varphi\)，并且
\[
F\cap \ker\varphi=0,
\]
因为若 \(x\in F\cap\ker\varphi\)，则 \(\varphi(x)=0\)，而 \(\varphi|_F\) 单射，所以 \(x=0\)。从而
\[
M\cong F\oplus\ker\varphi.
\]

设 \(m_1,\dots,m_r\) 生成 \(M\)。把每个 \(m_j\) 唯一写成
\[
m_j=k_j+f_j,\qquad k_j\in\ker\varphi,\ f_j\in F.
\]
现取任意 \(k\in\ker\varphi\)。因 \(m_j\) 生成 \(M\)，可写
\[
k=\sum_{j=1}^r a_jm_j=\sum_{j=1}^r a_jk_j+\sum_{j=1}^r a_jf_j.
\]
左端与第一项都在 \(\ker\varphi\) 中，故第二项也在 \(\ker\varphi\) 中；但它又在 \(F\) 中，所以它属于 \(F\cap\ker\varphi=0\)。因此
\[
k=\sum_{j=1}^r a_jk_j.
\]
这说明 \(\ker\varphi\) 由 \(k_1,\dots,k_r\) 生成，从而有限生成。

\setcounter{enumi}{12}
\item 设 \(f:A\to B\) 为环同态，\(N\) 是 \(B\)-模。把 \(N\) 看成 \(A\)-模后，证明自然映射
\[
g:N\to B\otimes_A N,\qquad y\mapsto 1\otimes y
\]
单射，且其像是一个直和项。

\textbf{证明。}
定义
\[
p:B\otimes_A N\to N,\qquad b\otimes y\mapsto by.
\]
先说明它是良定义的。张量积关系为
\[
(ba)\otimes y=b\otimes ay \qquad (a\in A),
\]
而把 \(N\) 看成 \(A\)-模时，\(ay=f(a)y\)，所以
\[
p((ba)\otimes y)=(ba)y=b(ay)=p(b\otimes ay),
\]
故 \(p\) 良定义。它显然是 \(B\)-线性的。

再看复合 \(p\circ g\)。对任意 \(y\in N\)，
\[
(p\circ g)(y)=p(1\otimes y)=1\cdot y=y.
\]
故
\[
p\circ g=\mathrm{id}_N.
\]
于是 \(g\) 有左逆，因此 \(g\) 必单射。

更进一步，因为 \(g\) 有左逆，短正合列
\[
0\to N\xrightarrow{g} B\otimes_A N\xrightarrow{} \operatorname{coker}(g)\to 0
\]
分裂，所以
\[
B\otimes_A N\cong g(N)\oplus \ker p.
\]
换言之，\(g(N)\) 是 \(B\otimes_A N\) 的一个直和项。
\end{enumerate}

\section{第3章：Rings and Modules of Fractions}

\begin{enumerate}[leftmargin=*, label=\textbf{3.\arabic*.}]
\setcounter{enumi}{1}
\item 设 \(\mathfrak a\) 是环 \(A\) 的一个理想，令 \(S=1+\mathfrak a\)。证明
\[
S^{-1}\mathfrak a \subseteq J(S^{-1}A).
\]
并据此配合 Nakayama 引理，给出命题 \((2.5)\) 的一个不依赖行列式的证明。

\textbf{证明。}
任取 \(x=s^{-1}a\in S^{-1}\mathfrak a\)，其中 \(a\in\mathfrak a,\ s\in S\)。则
\[
1-x=\frac{s-a}{s}.
\]
由于 \(s\in 1+\mathfrak a\)，可写成 \(s=1+a'\)（\(a'\in\mathfrak a\)），于是
\[
s-a=1+(a'-a)\in 1+\mathfrak a=S.
\]
故 \(1-x\) 在 \(S^{-1}A\) 中可逆。于是每个 \(x\in S^{-1}\mathfrak a\) 都属于 Jacobson 根，即
\[
S^{-1}\mathfrak a \subseteq J(S^{-1}A).
\]

现在设 \(M=\mathfrak aM\)。局部化后得
\[
S^{-1}M=(S^{-1}\mathfrak a)(S^{-1}M).
\]
由刚证结论与 Nakayama 引理，得到
\[
S^{-1}M=0.
\]
再由 Exercise 1，存在某个 \(s\in S\) 使
\[
sM=0.
\]
而 \(S=1+\mathfrak a\)，故存在 \(a\in\mathfrak a\) 使 \(s=1+a\)。于是
\[
(1+a)M=0.
\]
这正是 \((2.5)\) 的无行列式形式。

\item 设 \(A\) 是环，\(S,T\subseteq A\) 都是乘法闭集，记 \(U\) 为 \(T\) 在 \(S^{-1}A\) 中的像。证明
\[
(ST)^{-1}A \cong U^{-1}(S^{-1}A).
\]

\textbf{证明。}
两边都满足同一泛性质：它们都是把 \(A\) 中 \(S\cup T\) 的元素都变成可逆元所得到的局部化。

具体地，定义
\[
\phi:(ST)^{-1}A\to U^{-1}(S^{-1}A),\qquad \frac{a}{st}\mapsto \frac{a/s}{t/1},
\]
以及
\[
\psi:U^{-1}(S^{-1}A)\to (ST)^{-1}A,\qquad \frac{a/s}{t/1}\mapsto \frac{a}{st}.
\]
可以逐步核对这两个映射都良定义，且互为逆映射，因此得到所求同构。

\setcounter{enumi}{14}
\item 设 \(f:A\to B\) 是环同态。证明：
\begin{enumerate}[label=(\roman*)]
\item 若 \(S\subseteq A\) 乘法闭，则 \(\operatorname{Spec}(S^{-1}A)\) 自然同胚于其在 \(\operatorname{Spec}(A)\) 中的像；
\item 局部化后的谱映射是原谱映射在相应子空间上的限制；
\item 模掉理想后的谱映射也是原谱映射的限制；
\item 对 \(\mathfrak p\in\operatorname{Spec}(A)\)，纤维 \(f^{*-1}(\mathfrak p)\) 自然同胚于
\[
\operatorname{Spec}(B_{\mathfrak p}/\mathfrak p B_{\mathfrak p})
\cong
\operatorname{Spec}(k(\mathfrak p)\otimes_A B).
\]
\end{enumerate}

\textbf{证明。}
\begin{enumerate}[label=(\roman*)]
\item 由局部化的素理想对应，
\[
\operatorname{Spec}(S^{-1}A)\cong \{\mathfrak q\in \operatorname{Spec}(A):\mathfrak q\cap S=\varnothing\},
\]
这正是它在 \(\operatorname{Spec}(A)\) 中的像。

\item 在上述识别下，局部化后的谱映射显然就是原谱映射的限制。

\item 若 \(\mathfrak a\subseteq A\)，记 \(\mathfrak b=\mathfrak a^e\subseteq B\)，则
\[
\operatorname{Spec}(A/\mathfrak a)\cong V(\mathfrak a),\qquad
\operatorname{Spec}(B/\mathfrak b)\cong V(\mathfrak b),
\]
而诱导映射正是 \(f^*\) 在 \(V(\mathfrak b)\) 上的限制。

\item 取 \(S=A\setminus\mathfrak p\)。先在 \(\mathfrak p\) 处局部化，再对 \(S^{-1}\mathfrak p\) 取商，就得到纤维环
\[
B_{\mathfrak p}/\mathfrak p B_{\mathfrak p}.
\]
而
\[
B_{\mathfrak p}/\mathfrak p B_{\mathfrak p}\cong k(\mathfrak p)\otimes_A B,
\]
故得第二个同构。
\end{enumerate}

\setcounter{enumi}{18}
\item 设 \(M\) 是 \(A\)-模，定义
\[
\operatorname{Supp}(M)=\{\mathfrak p\in\operatorname{Spec}(A):M_{\mathfrak p}\neq 0\}.
\]
证明：
\begin{enumerate}[label=(\roman*)]
\item \(M\neq 0 \iff \operatorname{Supp}(M)\neq \varnothing\)；
\item \(\operatorname{Supp}(A/\mathfrak a)=V(\mathfrak a)\)；
\item 若 \(0\to M'\to M\to M''\to 0\) 正合，则
\[
\operatorname{Supp}(M)=\operatorname{Supp}(M')\cup \operatorname{Supp}(M'');
\]
\item \(\operatorname{Supp}(\bigoplus_i M_i)=\bigcup_i \operatorname{Supp}(M_i)\)；
\item 若 \(M\) 有限生成，则
\[
\operatorname{Supp}(M)=V(\operatorname{Ann}(M));
\]
\item 若 \(M,N\) 有限生成，则
\[
\operatorname{Supp}(M\otimes_A N)=\operatorname{Supp}(M)\cap \operatorname{Supp}(N);
\]
\item 若 \(M\) 有限生成，则
\[
\operatorname{Supp}(M/\mathfrak aM)=V(\mathfrak a+\operatorname{Ann}(M)).
\]
\end{enumerate}

\textbf{证明。}
\begin{enumerate}[label=(\roman*)]
\item 若 \(M\neq 0\)，取 \(0\neq x\in M\)，再取包含 \(\operatorname{Ann}(x)\) 的极大理想 \(\mathfrak m\)，则 \(x/1\neq 0\in M_{\mathfrak m}\)。反向显然。

\item
\[
(A/\mathfrak a)_{\mathfrak p}=0 \iff \mathfrak a\not\subseteq \mathfrak p,
\]
故支撑恰为 \(V(\mathfrak a)\)。

\item 局部化保持正合，因此逐点判断即可。

\item 因为局部化与直和可交换。

\item 由 Proposition 3.14，
\[
\operatorname{Ann}(M_{\mathfrak p})=A_{\mathfrak p}\operatorname{Ann}(M),
\]
故
\[
M_{\mathfrak p}=0 \iff \operatorname{Ann}(M)\not\subseteq \mathfrak p.
\]

\item 对每个 \(\mathfrak p\),
\[
(M\otimes_A N)_{\mathfrak p}\cong M_{\mathfrak p}\otimes_{A_{\mathfrak p}}N_{\mathfrak p}.
\]
若两者都非零且有限生成，则在局部环 \(A_{\mathfrak p}\) 上张量积仍非零。

\item 由
\[
M/\mathfrak aM \cong (A/\mathfrak a)\otimes_A M
\]
和上一条得到
\[
\operatorname{Supp}(M/\mathfrak aM)=V(\mathfrak a)\cap V(\operatorname{Ann}(M))
=V(\mathfrak a+\operatorname{Ann}(M)).
\]
\end{enumerate}

\setcounter{enumi}{20}
\item 设 \(S\subseteq A\) 乘法闭。证明 \(\operatorname{Spec}(S^{-1}A)\) 自然同胚于
\[
\{\mathfrak p\in \operatorname{Spec}(A):\mathfrak p\cap S=\varnothing\}.
\]

\textbf{证明。}
局部化的素理想恰好与不交于 \(S\) 的素理想对应：
\[
\mathfrak q\longmapsto \mathfrak q^c,\qquad
\mathfrak p\longmapsto S^{-1}\mathfrak p.
\]
这两个对应互为逆，并与拓扑兼容，因此给出同胚。

\setcounter{enumi}{21}
\item 设 \(\mathfrak p\in\operatorname{Spec}(A)\)。证明 \(\operatorname{Spec}(A_{\mathfrak p})\) 在 \(\operatorname{Spec}(A)\) 中的像等于所有包含 \(\mathfrak p\) 的开邻域之交。

\textbf{证明。}
\(\operatorname{Spec}(A_{\mathfrak p})\) 的像是
\[
\{\mathfrak q\in\operatorname{Spec}(A):\mathfrak q\subseteq \mathfrak p\}.
\]
若 \(\mathfrak q\subseteq \mathfrak p\)，则任一包含 \(\mathfrak p\) 的基本开邻域 \(D(f)\) 都满足 \(f\notin\mathfrak q\)，故 \(\mathfrak q\in D(f)\)。反之若 \(\mathfrak q\not\subseteq \mathfrak p\)，取 \(f\in\mathfrak q\setminus\mathfrak p\)，则 \(D(f)\) 含 \(\mathfrak p\) 但不含 \(\mathfrak q\)。故结论成立。
\end{enumerate}

\section{第4章：Primary Decomposition}

\begin{enumerate}[leftmargin=*, label=\textbf{4.\arabic*.}]
\setcounter{enumi}{1}
\item 若理想 \(\mathfrak a\) 满足 \(\mathfrak a=\sqrt{\mathfrak a}\)，证明它没有 embedded prime ideals。

\textbf{证明。}
设 \(A\) Noetherian，\(\mathfrak a\) 有一个不可冗余 primary decomposition
\[
\mathfrak a=\bigcap_{i=1}^r q_i,\qquad \sqrt{q_i}=\mathfrak p_i.
\]
由 \textit{Algebra III} 的 (III.Cor 6.2.12)，这些 \(\mathfrak p_i\) 恰是 \(A/\mathfrak a\) 的 associated primes。embedded prime 指其中不是极小素理想的那些。

现在因为 \(\mathfrak a\) 本身是根理想，故
\[
\mathfrak a=\sqrt{\mathfrak a}=\bigcap_{\mathfrak p\in \operatorname{Min}(A/\mathfrak a)}\mathfrak p
\]
即只用极小素理想就已经能恢复 \(\mathfrak a\)。因此在不可冗余 primary decomposition 中，不需要额外的非极小根理想成分。换言之，associated primes 全都已经是极小素理想，所以没有 embedded primes。

\setcounter{enumi}{3}
\item 在 \(\mathbb Z[t]\) 中，\(\mathfrak m=(2,t)\) 是极大理想，而 \(\mathfrak q=(4,t)\) 是 \(\mathfrak m\)-primary，但不是 \(\mathfrak m\) 的幂。证明之。

\textbf{证明。}
先证 \(\mathfrak m\) 极大。因为
\[
\mathbb Z[t]/(2,t)\cong (\mathbb Z/2\mathbb Z)[t]/(t)\cong \mathbb F_2
\]
是域，所以 \((2,t)\) 极大。

再证 \(\mathfrak q=(4,t)\) 是 \(\mathfrak m\)-primary。其根理想为
\[
\sqrt{(4,t)}=(2,t)=\mathfrak m.
\]
又商环
\[
\mathbb Z[t]/(4,t)\cong \mathbb Z/4\mathbb Z
\]
的唯一素理想是 \((2)/(4)\)，故其 nilradical 就是唯一极大理想。于是 \((4,t)\) 的根理想为极大理想，故 \((4,t)\) 是 \(\mathfrak m\)-primary。

最后说明它不是 \(\mathfrak m\) 的幂。注意
\[
\mathfrak m^2=(4,2t,t^2).
\]
对任意 \(n\ge 2\)，都有 \(\mathfrak m^n\subseteq \mathfrak m^2\)。但 \(t\in (4,t)\)，而 \(t\notin \mathfrak m^2\)，因为 \(\mathfrak m^2\) 中每项要么含系数 \(4\)，要么至少含两个 \((2,t)\) 因子，不可能得到单独的 \(t\)。因此
\[
(4,t)\neq \mathfrak m^n\qquad (n\ge 2).
\]
另一方面 \((4,t)\neq \mathfrak m\)，因为 \(2\in \mathfrak m\) 但 \(2\notin (4,t)\)。所以 \((4,t)\) 不是 \(\mathfrak m\) 的任何幂。

\setcounter{enumi}{4}
\item 在 \(K[x,y,z]\) 中令
\[
\mathfrak p_1=(x,y),\quad \mathfrak p_2=(x,z),\quad \mathfrak m=(x,y,z),\quad \mathfrak a=\mathfrak p_1\mathfrak p_2.
\]
证明
\[
\mathfrak a=\mathfrak p_1\cap \mathfrak p_2\cap \mathfrak m^2
\]
是一个 reduced primary decomposition，并判断 isolated 与 embedded 成分。

\textbf{证明。}
先直接计算
\[
\mathfrak a=(x,y)(x,z)=(x^2,xy,xz,yz).
\]

另一方面，
\[
\mathfrak p_1\cap\mathfrak p_2=(x,y)\cap(x,z)=(x,yz).
\]
再与
\[
\mathfrak m^2=(x^2,xy,xz,y^2,yz,z^2)
\]
相交。由于 \((x,yz)\) 中的元素写成 \(xu+yz\,v\)，要落入 \(\mathfrak m^2\) 就等价于常数项与一次项全部消失，于是正好得到
\[
(x,yz)\cap \mathfrak m^2=(x^2,xy,xz,yz)=\mathfrak a.
\]
故等式成立。

接着看 primary 性。\(\mathfrak p_1,\mathfrak p_2\) 是素理想，所以当然分别是 \(\mathfrak p_1\)-primary 与 \(\mathfrak p_2\)-primary。\(\mathfrak m\) 是极大理想，而 \(\mathfrak m^2\) 的根理想是 \(\mathfrak m\)，故 \(\mathfrak m^2\) 是 \(\mathfrak m\)-primary。

三个根理想
\[
\mathfrak p_1,\ \mathfrak p_2,\ \mathfrak m
\]
两两不同，因此这是 reduced decomposition。

最后判断 isolated / embedded。极小素理想是 \(\mathfrak p_1,\mathfrak p_2\)；因为
\[
\mathfrak a=\mathfrak p_1\mathfrak p_2\subseteq \mathfrak p_i
\]
且二者互不包含，所以它们是 \(V(\mathfrak a)\) 的两个极小点。故对应的 \(\mathfrak p_1,\mathfrak p_2\)-primary 成分是 isolated components。\(\mathfrak m\) 包含 \(\mathfrak p_1,\mathfrak p_2\)，不是极小素理想，因此 \(\mathfrak m^2\) 对应 embedded component。
\end{enumerate}

\section{第5章：Integral Dependence and Valuations}

\begin{enumerate}[leftmargin=*, label=\textbf{5.\arabic*.}]
\item 设 \(A\subset B\) 且 \(B\) 整于 \(A\)。若 \(f:A\to \Omega\) 映入代数闭域，证明 \(f\) 可延拓到 \(B\to \Omega\)。

\textbf{证明。}
设 \(\mathfrak p=\ker f\)。则 \(f\) 诱导单射
\[
\bar f:A/\mathfrak p \hookrightarrow \Omega.
\]
由整扩张上的 lying-over，即 (III.Thm 4.1.14)，存在 \(\mathfrak q\in \operatorname{Spec}B\) 使
\[
\mathfrak q\cap A=\mathfrak p.
\]
于是有整扩张
\[
A/\mathfrak p \subset B/\mathfrak q.
\]
而 \(B/\mathfrak q\) 是整域，因为 \(\mathfrak q\) 为素理想。令 \(L=\operatorname{Frac}(B/\mathfrak q)\)。由于 \(B/\mathfrak q\) 整于 \(A/\mathfrak p\)，所以 \(L\) 是 \(\operatorname{Frac}(A/\mathfrak p)\) 的代数扩张。

因 \(\Omega\) 代数闭，\(\bar f\) 可从 \(A/\mathfrak p\) 的分式域延拓到 \(L\)，特别可限制到 \(B/\mathfrak q\)。最后与商映射 \(B\to B/\mathfrak q\) 复合，就得到
\[
B\to \Omega
\]
且它在 \(A\) 上与 \(f\) 一致。

\setcounter{enumi}{3}
\item 若 \(B\) 整于 \(A\)，\(n\subset B\) 极大，\(m=n\cap A\)。问 \(B_n\) 是否总整于 \(A_m\)？

\textbf{证明。}
一般不成立。下面给出反例。

取
\[
A=k[x^2-1]\subset B=k[x],\qquad n=(x-1)\subset B.
\]
因为 \(x\) 满足首一方程
\[
T^2-(x^2-1)-1=0
\]
故 \(x\) 整于 \(A\)，所以 \(B=A[x]\) 整于 \(A\)。

设
\[
m=n\cap A=(x^2-1).
\]
则
\[
B_n=k[x]_{(x-1)},\qquad A_m=k[x^2-1]_{(x^2-1)}.
\]
现假设 \(\operatorname{char}k\neq 2\)，并设
\[
u=(x+1)^{-1}\in B_n.
\]
由
\[
(x^2-1)u=(x-1)
\]
两边再乘以 \(u\)，得到
\[
(x^2-1)u^2+2u-1=0.
\]
因此 \(u\) 在
\[
K=\operatorname{Frac}(A_m)=k(x^2-1)
\]
上至多是二次代数元。若 \(u\) 整于 \(A_m\)，则其范数
\[
N_{K(u)/K}(u)=-\frac{1}{x^2-1}
\]
也整于 \(A_m\)。但它已经属于分式域 \(K\)，而 \(A_m\) 在 \(K\) 中整闭，所以它应属于 \(A_m\)。这不可能，因为 \((x^2-1)^{-1}\notin A_m\)。

故 \(u\) 不整于 \(A_m\)，从而 \(B_n\) 不一定整于 \(A_m\)。

\setcounter{enumi}{4}
\item 设 \(B\) 整于 \(A\)。证明：
\begin{enumerate}[label=(\roman*)]
\item 若 \(x\in A\) 在 \(B\) 中是单位，则它在 \(A\) 中已是单位。
\item \(J(A)=J(B)\cap A\)。
\end{enumerate}

\textbf{证明。}
\begin{enumerate}[label=(\roman*)]
\item 设 \(x^{-1}\in B\)。因为 \(B\) 整于 \(A\)，故 \(x^{-1}\) 整于 \(A\)。于是存在首一方程
\[
(x^{-1})^n+a_1(x^{-1})^{n-1}+\cdots+a_n=0,\qquad a_i\in A.
\]
两边乘以 \(x^{n-1}\)，得
\[
x^{-1}+a_1+a_2x+\cdots+a_nx^{n-1}=0,
\]
即
\[
x^{-1}=-(a_1+a_2x+\cdots+a_nx^{n-1})\in A.
\]
所以 \(x\) 在 \(A\) 中可逆。

\item 先证 \(J(B)\cap A\subseteq J(A)\)。取 \(a\in J(B)\cap A\)。要证 \(a\) 落在 \(A\) 的每个极大理想中。任取极大理想 \(m\subset A\)。由 lying-over，即 (III.Thm 4.1.14)，存在极大理想 \(n\subset B\) 使 \(n\cap A=m\)；极大性可用 (III.Cor 4.1.12)。由于 \(a\in J(B)\subseteq n\)，故 \(a\in n\cap A=m\)。所以 \(a\in J(A)\)。

再证 \(J(A)\subseteq J(B)\cap A\)。取 \(a\in J(A)\)。任取 \(B\) 的极大理想 \(n\)，则 \(m=n\cap A\) 由 (III.Cor 4.1.12) 为极大理想。因 \(a\in J(A)\subseteq m\subseteq n\)，故 \(a\) 落在所有极大理想 \(n\) 中，即 \(a\in J(B)\)。

综上，
\[
J(A)=J(B)\cap A.
\]
\end{enumerate}

\setcounter{enumi}{6}
\item 若 \(A\subset B\) 且 \(B\setminus A\) 对乘法封闭，证明 \(A\) 在 \(B\) 中整闭。

\textbf{证明。}
设 \(x\in B\) 整于 \(A\)。需证 \(x\in A\)。

若 \(x\notin A\)，在一切满足首一整关系的多项式中取次数最小者，于是
\[
x^n+a_1x^{n-1}+\cdots+a_n=0,\qquad a_i\in A.
\]
由最小性，
\[
y:=x^{n-1}+a_1x^{n-2}+\cdots+a_{n-1}\notin A,
\]
否则 \(x\) 会满足次数更低的首一方程。又显然 \(x\notin A\)。由于 \(B\setminus A\) 对乘法封闭，故
\[
xy\notin A.
\]
但由整关系
\[
xy=-a_n\in A,
\]
矛盾。

故任何整于 \(A\) 的元素都已在 \(A\) 中，即 \(A\) 在 \(B\) 中整闭。

\setcounter{enumi}{11}
\item 设 \(G\) 是 \(A\) 的有限自同构群，\(A^G\) 为不动子环。证明 \(A\) 整于 \(A^G\)。

\textbf{证明。}
任取 \(x\in A\)，考虑多项式
\[
P_x(T)=\prod_{\sigma\in G}(T-\sigma(x)).
\]
它是首一多项式。其各系数是 \(\{\sigma(x)\}_{\sigma\in G}\) 的对称多项式，因此对任何 \(\tau\in G\) 都不变；也就是说这些系数都属于 \(A^G\)。又因为恒等元属于 \(G\)，代入 \(T=x\) 得
\[
P_x(x)=0.
\]
故 \(x\) 满足一个系数在 \(A^G\) 中的首一方程，所以 \(x\) 整于 \(A^G\)。

\setcounter{enumi}{12}
\item 在上题情形下，设 \(\mathfrak p\in \operatorname{Spec}(A^G)\)，\(P\) 为 \(A\) 中所有收缩为 \(\mathfrak p\) 的素理想集合。证明 \(G\) 在 \(P\) 上传递。

\textbf{证明。}
先说明 \(G\) 的作用保持 \(P\)。若 \(\mathfrak q\cap A^G=\mathfrak p\)，则对任意 \(\sigma\in G\)，因为 \(\sigma\) 固定 \(A^G\) 上的每个元素，故
\[
\sigma(\mathfrak q)\cap A^G=\mathfrak p.
\]
所以 \(\sigma(\mathfrak q)\in P\)。

现取 \(\mathfrak q_1,\mathfrak q_2\in P\)。我们证明存在 \(\sigma\in G\) 使
\[
\sigma(\mathfrak q_2)=\mathfrak q_1.
\]
反设不存在。则所有 \(\sigma(\mathfrak q_2)\) 都不包含 \(\mathfrak q_1\)。由 prime avoidance，可取
\[
x\in \mathfrak q_1\setminus \bigcup_{\sigma\in G}\sigma(\mathfrak q_2).
\]
于是对每个 \(\sigma\)，有 \(\sigma(x)\notin \mathfrak q_2\)。

但由上题构造的
\[
N(x)=\prod_{\sigma\in G}\sigma(x)
\]
属于 \(A^G\)。又因为 \(x\in\mathfrak q_1\)，所以 \(N(x)\in \mathfrak q_1\cap A^G=\mathfrak p\subseteq \mathfrak q_2\)。由于 \(\mathfrak q_2\) 是素理想，只要乘积在 \(\mathfrak q_2\) 中，就必有某个因子 \(\sigma(x)\in \mathfrak q_2\)，矛盾。

故作用必传递。

\setcounter{enumi}{13}
\item 设 \(A\) 整闭，分式域为 \(K\)，\(L/K\) 为有限正规可分扩张，\(B\) 是 \(A\) 在 \(L\) 中的整闭包，\(G=\operatorname{Gal}(L/K)\)。证明 \(\sigma(B)=B\) 且 \(A=B^G\)。

\textbf{证明。}
先证 \(\sigma(B)=B\)。任取 \(b\in B\)，则 \(b\) 整于 \(A\)，故有
\[
b^n+a_1b^{n-1}+\cdots+a_n=0,\qquad a_i\in A.
\]
对等式施以 \(\sigma\in G\)，注意 \(\sigma\) 在 \(K\) 上恒等，而 \(A\subset K\)，故 \(\sigma(a_i)=a_i\)。于是
\[
\sigma(b)^n+a_1\sigma(b)^{n-1}+\cdots+a_n=0.
\]
这表明 \(\sigma(b)\) 也整于 \(A\)，因此 \(\sigma(b)\in B\)。所以 \(\sigma(B)\subseteq B\)。对 \(\sigma^{-1}\) 同理，得 \(B\subseteq \sigma(B)\)，故 \(\sigma(B)=B\)。

再证 \(A=B^G\)。显然 \(A\subseteq B^G\)，因为 \(A\subset K=L^G\)，且 \(A\subset B\)。

反过来，若 \(x\in B^G\)，则 \(x\in L^G=K\)。另一方面 \(x\in B\)，故 \(x\) 整于 \(A\)。因为 \(A\) 在其分式域 \(K\) 中整闭，所以 \(x\in A\)。故 \(B^G\subseteq A\)。

综上，
\[
B^G=A.
\]

\setcounter{enumi}{14}
\item 设 \(A\) 如上，\(L/K\) 有限扩张，\(B\) 为 \(A\) 在 \(L\) 中的整闭包。证明 \(\operatorname{Spec}(B)\to \operatorname{Spec}(A)\) 的纤维有限。

\textbf{证明。}
先处理 \(L/K\) 为有限 Galois 扩张的情形。设 \(G=\operatorname{Gal}(L/K)\)。对给定 \(\mathfrak p\in\operatorname{Spec}A\)，所有收缩到 \(\mathfrak p\) 的素理想集合在上题中被证明是一个 \(G\)-轨道，因此其基数至多为 \(|G|\)，从而有限。

一般有限可分扩张情形，可嵌入某个有限 Galois 扩张 \(L'/K\)。设 \(B'\) 为 \(A\) 在 \(L'\) 中的整闭包，则 \(B'\) 整于 \(B\)。于是 \(\operatorname{Spec}(B')\to \operatorname{Spec}(B)\) 满射，由 (III.Thm 4.1.14)。若 \(B'\) 上某条纤维有限，则 \(B\) 上对应纤维作为其像，也有限。

对纯不可分扩张，可用唯一性。若 \(x\in L\)，则存在 \(p^n\) 使 \(x^{p^n}\in K\)。对于收缩到 \(\mathfrak p\) 的素理想，元素是否落入该素理想由某个 \(p^n\) 次幂是否落入 \(\mathfrak p\) 决定，从而纤维至多一个点。于是也有限。

把一般有限扩张分解为先纯不可分、后可分（或反之）的复合，即得总纤维有限。
\end{enumerate}

\section{第6章：Chain Conditions}

\begin{enumerate}[leftmargin=*, label=\textbf{6.\arabic*.}]
\item 若 \(M\) Noetherian 且 \(u:M\to M\) 满射，证明 \(u\) 为同构；若 \(M\) Artinian 且 \(u\) 单射，也证明 \(u\) 为同构。

\textbf{证明。}
先证 Noetherian 情形。考虑核的升链
\[
\ker u\subseteq \ker u^2\subseteq \ker u^3\subseteq \cdots.
\]
因 \(M\) Noetherian，此链稳定：存在 \(n\) 使
\[
\ker u^n=\ker u^{n+1}.
\]
记这一稳定的核为 \(K\)。由满射性，
\[
u(K)=K:
\]
若 \(z\in K\)，取 \(y\) 使 \(u(y)=z\)，则
\[
u^{n+1}(y)=u^n(z)=0,
\]
所以 \(y\in \ker u^{n+1}=K\)，从而 \(z=u(y)\in u(K)\)。

现在若 \(x\in K\)，则因 \(K=u(K)\)，可反复写成
\[
x=u(x_1)=u^2(x_2)=\cdots=u^n(x_n)
\]
其中每个 \(x_i\in K\)。于是
\[
u^n(x_n)=x,\qquad x_n\in K=\ker u^n,
\]
故 \(x=0\)。所以 \(K=0\)，特别 \(\ker u=0\)。因此 \(u\) 既满又单，是同构。

再证 Artinian 情形。考虑像的降链
\[
M\supseteq u(M)\supseteq u^2(M)\supseteq \cdots.
\]
因 \(M\) Artinian，此链稳定：存在 \(n\) 使
\[
I:=u^n(M)=u^{n+1}(M).
\]
于是 \(u\) 在商模 \(M/I\) 上诱导单射
\[
\bar u:M/I\to M/I.
\]
并且 \(\bar u\) 是幂零的，因为
\[
\bar u^n(M/I)=u^n(M)/I=I/I=0.
\]
一个单射不可能是非零模上的幂零自同态，所以 \(M/I=0\)。因此 \(I=M\)，即 \(u^n(M)=M\)，从而 \(u(M)=M\)。故 \(u\) 也满射，于是是同构。

\item 若 \(M\) 的任意非空有限生成子模集合都有极大元，证明 \(M\) Noetherian。

\textbf{证明。}
要证 \(M\) Noetherian，只需证每个子模都有限生成，这正是 \textit{Algebra III} 中 (III.Prop 5.1.3) 的判别。

取任意子模 \(N\subseteq M\)。若 \(N=0\) 则显然有限生成。若 \(N\neq 0\)，考虑所有 \(N\) 的有限生成子模所成集合
\[
\mathcal F=\{L\subseteq N: L \text{ 有限生成}\}.
\]
它非空，因为任取 \(0\neq x\in N\)，则 \(Ax\in\mathcal F\)。依假设，\(\mathcal F\) 有极大元，记为 \(N_0\)。

若 \(N_0\neq N\)，取 \(x\in N\setminus N_0\)。则
\[
N_0+Ax
\]
仍是有限生成子模，且严格包含 \(N_0\)，这与 \(N_0\) 的极大性矛盾。故 \(N=N_0\)，从而 \(N\) 有限生成。

任意子模都有限生成，所以 \(M\) Noetherian。
\end{enumerate}

\section{第7章：Noetherian Rings}

\begin{enumerate}[leftmargin=*, label=\textbf{7.\arabic*.}]
\item 设 \(A\) 是 Dedekind 域，\(S\subset A\) 乘法闭。证明 \(S^{-1}A\) 不是 Dedekind 域就是 \(A\) 的分式域；并证明类群映射满射。

\textbf{证明。}
Dedekind 域按 \textit{Algebra III} 的刻画可用 (III.Thm 7.2.9)：每个非零分式理想都可逆；等价地，每个局部化于非零素理想处都是 DVR。

对 \(S^{-1}A\) 而言，若 \(S\) 与 \(A\) 的每个非零素理想都相交，则所有非零素理想在局部化后都消失，故 \(S^{-1}A\) 没有非零素理想，只能是分式域 \(\operatorname{Frac}(A)\)。

否则，存在非零素理想 \(\mathfrak p\) 与 \(S\) 不相交。由局部化与谱的对应，\(S^{-1}A\) 的非零素理想正来自这些与 \(S\) 不相交的非零素理想。对这样的 \(\mathfrak p\)，有
\[
(S^{-1}A)_{S^{-1}\mathfrak p}\cong A_{\mathfrak p}.
\]
而 \(A\) 是 Dedekind 域，所以 \(A_{\mathfrak p}\) 是 DVR，即 (III.Cor 7.1.5) 与 (III.Prop 7.1.4)。故 \(S^{-1}A\) 在每个非零素理想处局部化仍是 DVR，因而仍是 Dedekind 域。

类群满射也可从局部化看出。\(S^{-1}A\) 的任一非零分式理想都可写成 \(S^{-1}I\) 的形式，其中 \(I\) 是 \(A\) 的分式理想。于是
\[
\operatorname{Cl}(A)\to \operatorname{Cl}(S^{-1}A),\qquad [I]\mapsto [S^{-1}I]
\]
是满射。

\item 设 \(A\) 是 Dedekind 域，证明 Gauss 引理
\[
c(fg)=c(f)c(g).
\]

\textbf{证明。}
记 \(c(f)\) 为多项式 \(f\) 的 content 理想。显然总有
\[
c(fg)\subseteq c(f)c(g),
\]
因为 \(fg\) 的每个系数都是 \(f\) 的系数与 \(g\) 的系数的有限和。

反向包含可在每个极大理想处局部化验证。任取极大理想 \(\mathfrak m\)。Dedekind 域局部化到 \(\mathfrak m\) 处是 DVR，这在 \textit{Algebra III} 中是 (III.Cor 7.1.5) 或 (III.Prop 7.1.4) 的内容。设其统一参数元为 \(\pi\)。于是
\[
c(f)A_{\mathfrak m}=(\pi^r),\qquad c(g)A_{\mathfrak m}=(\pi^s).
\]
可写
\[
f=\pi^r f_0,\qquad g=\pi^s g_0
\]
且 \(f_0,g_0\) 至少各有一个系数为单位，所以
\[
c(f_0)=c(g_0)=A_{\mathfrak m}.
\]
于是
\[
fg=\pi^{r+s}f_0g_0.
\]
因为 \(f_0g_0\) 仍有某个系数为单位，故
\[
c(fg)A_{\mathfrak m}=(\pi^{r+s})=c(f)A_{\mathfrak m}\,c(g)A_{\mathfrak m}.
\]
对所有极大理想 \(\mathfrak m\) 成立，于是全局上有
\[
c(fg)=c(f)c(g).
\]

\setcounter{enumi}{11}
\item 若 \(B\) 是 faithfully flat 的 \(A\)-代数，且 \(B\) Noetherian，证明 \(A\) Noetherian。

\textbf{证明。}
设
\[
\mathfrak a_1\subseteq \mathfrak a_2\subseteq \cdots
\]
是 \(A\) 中任意理想升链。对每项扩张到 \(B\) 中，得
\[
\mathfrak a_1B\subseteq \mathfrak a_2B\subseteq \cdots.
\]
因为 \(B\) Noetherian，升链稳定：存在 \(N\) 使得对所有 \(n\ge N\),
\[
\mathfrak a_nB=\mathfrak a_{n+1}B.
\]

现在利用 faithfully flat。若 \(I\subseteq J\) 是 \(A\) 的两个理想且 \(IB=JB\)，则
\[
(J/I)\otimes_A B\cong JB/IB=0.
\]
faithful flat 意味着有限生成模乃至任意模在张量后为零只可能原来就为零，因此 \(J/I=0\)，即 \(I=J\)。

应用于上面的升链，得
\[
\mathfrak a_n=\mathfrak a_{n+1}\qquad (n\ge N).
\]
所以 \(A\) 满足理想升链条件，即 \(A\) Noetherian。这里也可对照 \textit{Algebra III} 中 Noether 性判据 (III.Prop 5.1.3)。
\end{enumerate}

\section{第8章：Artin Rings}

\begin{enumerate}[leftmargin=*, label=\textbf{8.\arabic*.}]
\setcounter{enumi}{1}
\item 设 \(A\) 是 Noetherian 环，\(f=\sum_{n\ge 0}a_nx^n\in A[[x]]\)。证明
\[
f \text{ 幂零}
\Longleftrightarrow
\forall n\ge 0,\ a_n \text{ 幂零}.
\]

\textbf{证明。}
若 \(f\) 幂零，则对任意素理想 \(\mathfrak p\)，其像在 \((A/\mathfrak p)[[x]]\) 中也幂零。由于 \(A/\mathfrak p\) 是整环，\((A/\mathfrak p)[[x]]\) 也是整环，因此幂零元只能是零元，故每个 \(a_n\in \mathfrak p\)。这对所有 \(\mathfrak p\) 成立，所以每个 \(a_n\in \operatorname{Nil}(A)\)，于是每个 \(a_n\) 都幂零。

反过来，设所有 \(a_n\) 都幂零。它们生成的理想
\[
I=(a_0,a_1,\dots)
\]
在 Noetherian 性下有限生成，而有限个幂零元生成的理想是幂零理想，故 \(I^N=0\) 对某个 \(N\) 成立。因为 \(f\in I[[x]]\)，所以
\[
f^N\in I^N[[x]]=0.
\]
故 \(f\) 幂零。

\item 设 \(\mathfrak a\) 是环 \(A\) 中的不可约理想。证明下列条件等价：
\begin{enumerate}[label=(\roman*)]
\item \(\mathfrak a\) 是 primary ideal；
\item 对每个乘法闭集 \(S\subseteq A\)，都有
\[
(S^{-1}\mathfrak a)^c=(\mathfrak a:x)
\quad\text{对某个 }x\in S;
\]
\item 对每个 \(x\in A\)，链
\[
(\mathfrak a:x)\subseteq (\mathfrak a:x^2)\subseteq \cdots
\]
稳定。
\end{enumerate}

\textbf{证明。}
\((i)\Rightarrow(ii)\)\; 对任意乘法闭集 \(S\)，总有
\[
(S^{-1}\mathfrak a)^c=\bigcup_{s\in S}(\mathfrak a:s).
\]
若 \(\mathfrak a\) 是 \(\mathfrak p\)-primary，则分两种情形：

若 \(S\cap \mathfrak p\neq\varnothing\)，取 \(s\in S\cap \mathfrak p\)。因为 \(\mathfrak a\) primary，某个幂 \(s^m\in\mathfrak a\)，且 \(s^m\in S\)。于是
\[
(\mathfrak a:s^m)=A,
\]
从而 \((S^{-1}\mathfrak a)^c=A=(\mathfrak a:s^m)\)。

若 \(S\cap \mathfrak p=\varnothing\)，则对任意 \(s\in S\)，由 \(sy\in\mathfrak a\) 可推出 \(y\in\mathfrak a\)。因此
\[
(S^{-1}\mathfrak a)^c=\mathfrak a=(\mathfrak a:1).
\]
故 (ii) 成立。

\((ii)\Rightarrow(iii)\)\; 取 \(S=\{1,x,x^2,\dots\}\)。则
\[
(S^{-1}\mathfrak a)^c=\bigcup_{n\ge 1}(\mathfrak a:x^n).
\]
由 (ii) 知该并等于某个 \((\mathfrak a:x^m)\)，故链自 \(m\) 起稳定。

\((iii)\Rightarrow(i)\)\; 固定 \(x\in A\)，设链从 \(n\) 起稳定。则有恒等式
\[
\mathfrak a=(\mathfrak a+(x^n))\cap (\mathfrak a:x^n).
\]
的确，右包含左显然；反过来若
\[
y=a+bx^n\in (\mathfrak a:x^n),
\]
则 \(yx^n\in\mathfrak a\)，故 \(bx^{2n}\in\mathfrak a\)，即
\[
b\in (\mathfrak a:x^{2n})=(\mathfrak a:x^n),
\]
于是 \(bx^n\in\mathfrak a\)，所以 \(y\in\mathfrak a\)。

由于 \(\mathfrak a\) 不可约，只能有
\[
\mathfrak a=\mathfrak a+(x^n)\quad\text{或}\quad \mathfrak a=(\mathfrak a:x^n).
\]
现在若 \(yz\in\mathfrak a\) 且 \(y\notin\mathfrak a\)，则 \(y\in(\mathfrak a:z)\)。对 \(x=z\) 应用上面的结论：若 \((\mathfrak a:z^n)=\mathfrak a\)，则
\[
y\in (\mathfrak a:z)\subseteq (\mathfrak a:z^n)=\mathfrak a,
\]
矛盾。故只能有
\[
\mathfrak a=\mathfrak a+(z^n),
\]
即 \(z^n\in\mathfrak a\)。因此模 \(\mathfrak a\) 的每个零因子都是幂零元，故 \(\mathfrak a\) 是 primary。

\setcounter{enumi}{3}
\item 判断下列环哪些是 Noetherian：
\begin{enumerate}[label=(\roman*)]
\item 在圆周 \(|z|=1\) 上无极点的有理函数环；
\item 正收敛半径的幂级数环；
\item 无限收敛半径的幂级数环；
\item 在原点前 \(k\) 阶导数全为零的多项式环；
\item 对 \(w\) 的所有偏导在 \(z=0\) 都为零的 \(\mathbb C[z,w]\) 的子环。
\end{enumerate}

\textbf{证明。}
\begin{enumerate}[label=(\roman*)]
\item 是。它是 \(\mathbb C[z]\) 的局部化，局部化保持 Noetherian。
\item 是。它就是一元收敛幂级数环 \(\mathbb C\{z\}\)。任一非零元都可唯一写成
\[
z^m u(z),
\]
其中 \(u(0)\neq 0\)，故 \(u(z)\) 是单位。因此每个理想都形如 \((z^m)\)，所以它是 DVR，特别是 Noetherian。
\item 不是。它就是整函数环，而整函数环不是 Noetherian。
\item 是。该环同构于有限生成 \(\mathbb C\)-代数
\[
\mathbb C[z^{k+1},z^{k+2},\dots,z^{2k+1}],
\]
因此 Noetherian。
\item 是。条件等价于多项式可被 \(z\) 整除，因此该子环就是
\[
\mathbb C[z]+z\mathbb C[z,w]=\mathbb C[z,zw],
\]
是有限生成 \(\mathbb C\)-代数，所以 Noetherian。
\end{enumerate}
\end{enumerate}

\section{第9章：Discrete Valuation Rings and Dedekind Domains}

\begin{enumerate}[leftmargin=*, label=\textbf{9.\arabic*.}]
\setcounter{enumi}{1}
\item 设 \(A\) 是 Dedekind 域，证明若 \(f\) 的 content 记为 \(c(f)\)，则
\[
c(fg)=c(f)c(g).
\]

\textbf{证明。}
这与第7章习题 2 完全相同。

关键输入仍然是：对每个极大理想 \(\mathfrak m\)，局部化 \(A_{\mathfrak m}\) 是 DVR，
这正是 (III.Cor 7.1.5) 或 (III.Prop 7.1.4) 的内容。
随后在 DVR 中用统一参数元计算 content 的赋值，
便得到
\[
c(fg)A_{\mathfrak m}=c(f)A_{\mathfrak m}\,c(g)A_{\mathfrak m}
\]
对所有 \(\mathfrak m\) 成立，因此全局上
\[
c(fg)=c(f)c(g).
\]

\setcounter{enumi}{3}
\item 设 \(A\) 是局部整环，不是域，极大理想 \(\mathfrak m\) 主生成，并满足
\[
\bigcap_{n\ge 1}\mathfrak m^n=0.
\]
证明 \(A\) 是 DVR。

\textbf{证明。}
设 \(\mathfrak m=(\pi)\)。因为 \(A\) 不是域，\(\pi\) 不是单位。

对任意非零 \(x\in A\)，考虑所有满足 \(x\in \mathfrak m^n\) 的 \(n\)。由于 \(\bigcap_{n\ge 1}\mathfrak m^n=0\)，这个集合不是全体非负整数，因此存在最大整数 \(v(x)\ge 0\) 使
\[
x\in \mathfrak m^{v(x)}\setminus \mathfrak m^{v(x)+1}.
\]
因为 \(\mathfrak m=(\pi)\)，可写
\[
x=\pi^{v(x)}u.
\]
若 \(u\in \mathfrak m\)，则 \(x\in \mathfrak m^{v(x)+1}\)，与 \(v(x)\) 的定义矛盾，所以 \(u\notin \mathfrak m\)，即 \(u\) 是单位。于是每个非零元都唯一写成
\[
x=u\pi^n,\qquad u\in A^\times,\ n\ge 0.
\]

把 \(v\) 延拓到分式域 \(K=\operatorname{Frac}(A)\) 上：若 \(x=a/b\) 且 \(a,b\neq 0\)，定义
\[
v(x)=v(a)-v(b).
\]
这给出一个离散赋值，而
\[
A=\{x\in K:v(x)\ge 0\}.
\]
因此 \(A\) 是一个离散赋值环。与 \textit{Algebra III} 中 DVR 判别 (III.Prop 7.1.4) 相符。

\setcounter{enumi}{6}
\item 设 \(A\) 是 Dedekind 域，\(\mathfrak a\neq 0\)。证明 \(A/\mathfrak a\) 中每个理想都是主理想，并推出 \(A\) 中每个理想都至多由两个元素生成。

\textbf{证明。}
Dedekind 域中每个非零理想唯一分解为素理想幂乘积，这是 (III.Cor 7.1.7)。设
\[
\mathfrak a=\prod_i \mathfrak p_i^{r_i}.
\]
由中国剩余定理，即 \textit{Algebra III} 的 (III.Cor 5.3.15)，有
\[
A/\mathfrak a \cong \prod_i A/\mathfrak p_i^{r_i}.
\]
因此只需证明每个 \(A/\mathfrak p^r\) 中理想都主生成。

而 \(A_{\mathfrak p}\) 是 DVR，由 (III.Prop 7.1.4)。在 DVR 中，每个理想都是某个 \((\pi^m)\)。模去 \(\mathfrak p^r\) 后，\(A/\mathfrak p^r\) 的理想就对应于
\[
\mathfrak p^m/\mathfrak p^r\qquad (0\le m\le r),
\]
因此都由 \(\pi^m\) 的像生成。故 \(A/\mathfrak a\) 中每个理想都是主理想。

现取 \(A\) 中任意非零理想 \(\mathfrak b\)。取 \(0\neq a\in \mathfrak b\)。则 \(\mathfrak b/(a)\) 是 \(A/(a)\) 的理想，故主生成。设它由 \(\bar y\) 生成。于是 \(\mathfrak b\) 中每个元素模 \((a)\) 都是 \(\bar y\) 的倍数，即
\[
\mathfrak b=(a,y).
\]
故每个非零理想至多由两个元素生成。零理想显然由一个元素生成。

\setcounter{enumi}{8}
\item 证明 Dedekind 域上的中国剩余定理的兼容形式：同余组
\[
x\equiv x_i \pmod{\mathfrak a_i}
\]
可解，当且仅当对所有 \(i\ne j\),
\[
x_i\equiv x_j \pmod{\mathfrak a_i+\mathfrak a_j}.
\]

\textbf{证明。}
必要性显然：若存在 \(x\) 同时满足全部同余，则
\[
x_i-x_j=(x_i-x)+(x-x_j)\in \mathfrak a_i+\mathfrak a_j.
\]

下面证充分性。由 Dedekind 域中的理想分解与局部化，只需对每个非零素理想 \(\mathfrak p\) 分别验证。局部化到 \(A_{\mathfrak p}\) 后，它是 DVR，由 (III.Prop 7.1.4)，故所有理想都是
\[
(\pi^{n_i})
\]
的形式，并按包含关系全序排列。

因此，在某个固定 \(\mathfrak p\)-局部化中，一组同余条件
\[
x\equiv x_i \pmod{\pi^{n_i}}
\]
只需检查对最大指数的那一条是否与其余兼容；而题设的
\[
x_i\equiv x_j \pmod{\mathfrak a_i+\mathfrak a_j}
\]
局部化后正好保证这种兼容性。于是每个局部系统都可解。

再把各个素理想方向上的局部解用标准中国剩余拼回，即得全局解。
\end{enumerate}

\section{第10章：Completions}

\begin{enumerate}[leftmargin=*, label=\textbf{10.\arabic*.}]
\setcounter{enumi}{2}
\item 设 \(A\) Noetherian，\(\mathfrak a\subset A\) 为理想，\(M\) 有限生成。证明
\[
\bigcap_{n\ge 0}\mathfrak a^n M
=
\bigcap_{\mathfrak m\supseteq \mathfrak a}
\ker\!\bigl(M\to \widehat{M_{\mathfrak m}}\bigr),
\]
并推出 \(\widehat M=0\) 当且仅当 \(\operatorname{Supp}(M)\cap V(\mathfrak a)=\varnothing\)。

\textbf{证明。}
设
\[
N=\bigcap_{n\ge 0}\mathfrak a^nM.
\]
由 Krull 交定理在 \textit{Algebra III} 中的形式，即 (III.Cor 8.2.4)，知 \(N\) 满足
\[
\mathfrak aN=N.
\]

另一方面，由完成映射的核描述，即 (III.Cor 8.5.6)，有限生成模 \(M_{\mathfrak m}\) 到其 \(\mathfrak aA_{\mathfrak m}\)-进完成的核正是
\[
\bigcap_{n\ge 0}\mathfrak a^n M_{\mathfrak m}.
\]
而局部化与有限交相容，所以
\[
\bigcap_{n\ge 0}\mathfrak a^n M_{\mathfrak m}
=
\left(\bigcap_{n\ge 0}\mathfrak a^n M\right)_{\mathfrak m}
=N_{\mathfrak m}.
\]
因此
\[
x\in \ker(M\to \widehat{M_{\mathfrak m}})\iff x/1\in N_{\mathfrak m}.
\]

先证 \(N\) 包含于右端交。若 \(x\in N\)，则其在每个局部化 \(M_{\mathfrak m}\) 中都落入 \(N_{\mathfrak m}\)，故映到每个 \(\widehat{M_{\mathfrak m}}\) 中都为零。

反过来，若 \(x\in M\) 映到每个 \(\widehat{M_{\mathfrak m}}\) 中都为零，则 \(x/1\in N_{\mathfrak m}\) 对每个 \(\mathfrak m\supseteq \mathfrak a\) 成立，即 \((Ax+N)_{\mathfrak m}=0\) 对每个这样的 \(\mathfrak m\) 成立。于是
\[
\operatorname{Supp}(Ax+N)\cap V(\mathfrak a)=\varnothing.
\]
由支持的标准判别，这等价于存在 \(t\in \mathfrak a\) 使
\[
t(Ax+N)=Ax+N.
\]
再对有限生成模 \(Ax+N\) 用 Nakayama（或 Krull 交的 Jacobson 形式）可得 \(Ax+N=0\)，即 \(x\in N\)。故等式成立。

再证第二部分。由 (III.Cor 8.5.6)，
\[
\widehat M=0 \iff \bigcap_{n\ge 0}\mathfrak a^nM=M.
\]
由 Krull 交定理 (III.Cor 8.2.4)，这等价于
\[
\mathfrak a M=M.
\]
对有限生成模，\(\mathfrak aM=M\) 当且仅当对每个包含 \(\mathfrak a\) 的极大理想 \(\mathfrak m\)，有 \(M_{\mathfrak m}=0\)；这又等价于
\[
\operatorname{Supp}(M)\cap V(\mathfrak a)=\varnothing.
\]
故结论成立。

\setcounter{enumi}{3}
\item 设 \(A\) Noetherian，\(\widehat A\) 为 \(\mathfrak a\)-进完成。证明若 \(x\in A\) 在 \(A\) 中不是零因子，则它在 \(\widehat A\) 中也不是零因子。并说明“\(A\) 为整环 \(\Rightarrow \widehat A\) 为整环”一般不成立。

\textbf{证明。}
因为 \(x\) 在 \(A\) 中不是零因子，乘法映射
\[
0\to A\xrightarrow{\cdot x} A
\]
单射。把它补成短正合列
\[
0\to A\xrightarrow{\cdot x}A\to A/xA\to 0.
\]
对该短正合列取 \(\mathfrak a\)-进完成。由于 \(A\) Noetherian，有限生成模的完成保持正合，这正是 \textit{Algebra III} 的 (III.Prop 8.5.1)。故得到
\[
0\to \widehat A\xrightarrow{\cdot x}\widehat A\to \widehat{A/xA}\to 0.
\]
第一箭头仍单射，因此 \(x\) 在 \(\widehat A\) 中也不是零因子。

但这并不推出 \(\widehat A\) 一定是整环，因为完成后可能出现新的零因子，它们未必来自原环中的元素。经典例子是
\[
A=k[x,y]_{(x,y)}/(y^2-x^2-x^3).
\]
原环对应局部尖点曲线，仍是整环；但其完成为
\[
k[[x,y]]/(y^2-x^2-x^3).
\]
在 \(k[[x]]\) 中若 \(1+x\) 有平方根，则
\[
y^2-x^2(1+x)=(y-x\sqrt{1+x})(y+x\sqrt{1+x}),
\]
从而完成环分裂，不再是整环。故“整环完成仍整”一般不成立。

\setcounter{enumi}{5}
\item 设 \(A\) Noetherian，\(\mathfrak a\subset A\)。证明
\[
\mathfrak a\subseteq J(A)
\Longleftrightarrow
\text{每个极大理想在 }\mathfrak a\text{-进拓扑下都是闭的}.
\]

\textbf{证明。}
先设 \(\mathfrak a\subseteq J(A)\)。任取极大理想 \(\mathfrak m\)。则 \(\mathfrak a\subseteq \mathfrak m\)，故对任意 \(n\),
\[
\mathfrak m+\mathfrak a^n=\mathfrak m.
\]
于是 \(\mathfrak m\) 的 \(\mathfrak a\)-进闭包为
\[
\bigcap_{n\ge 0}(\mathfrak m+\mathfrak a^n)=\mathfrak m,
\]
故 \(\mathfrak m\) 闭。

反过来，若存在极大理想 \(\mathfrak m\) 不含 \(\mathfrak a\)，取
\[
t\in \mathfrak a\setminus \mathfrak m.
\]
则 \(t\) 在域 \(A/\mathfrak m\) 中是单位，所以对每个 \(n\),
\[
\mathfrak m+\mathfrak a^n=A.
\]
于是 \(\mathfrak m\) 的闭包为
\[
\bigcap_{n\ge 0}(\mathfrak m+\mathfrak a^n)=A,
\]
不是闭集。故若每个极大理想都闭，就必有 \(\mathfrak a\subseteq J(A)\)。

这与 \textit{Algebra III} 对 Zariski ring 的讨论一致，见 (III.Prop 8.5.4) 附近。

\setcounter{enumi}{6}
\item 设 \(\widehat A\) 是 \(A\) 的 \(\mathfrak a\)-进完成。证明
\[
\widehat A \text{ faithfully flat over } A
\Longleftrightarrow
A \text{ 是 Zariski ring}.
\]

\textbf{证明。}
在 Noetherian 情形，完成总是 flat，这是 \textit{Algebra III} 的 (III.Prop 8.5.4)。
因此只需判别何时 faithfully flat。

一方面，若 \(\widehat A\) faithfully flat over \(A\)，则对任意有限生成 \(A\)-模 \(M\neq 0\)，有
\[
M\otimes_A \widehat A=\widehat M\neq 0.
\]
于是 \(\widehat M=0\) 不会发生在非零 \(M\) 上。由上一题与前一题可知，这等价于不存在满足 \(\operatorname{Supp}(M)\cap V(\mathfrak a)=\varnothing\) 的所有局部行为失效情形；最终等价于
\[
\mathfrak a\subseteq J(A).
\]

另一方面，若 \(\mathfrak a\subseteq J(A)\)，则对任意有限生成 \(M\neq 0\)，由上一题知
\[
\operatorname{Supp}(M)\cap V(\mathfrak a)\neq \varnothing,
\]
所以 \(\widehat M\neq 0\)。这正是 faithful flatness 的模判据。

而 \(\mathfrak a\subseteq J(A)\) 正是 Zariski ring 的定义之一，所以结论成立。

\setcounter{enumi}{8}
\item 证明 Hensel 引理：若 \((A,\mathfrak m)\) 是 \(\mathfrak m\)-进完备局部环，\(f(x)\in A[x]\) 首一，且其模 \(\mathfrak m\) 的像分解为互素首一多项式
\[
\bar f=\bar g\,\bar h,
\]
则可提升为
\[
f=gh
\]
其中 \(g,h\in A[x]\) 仍首一，且模 \(\mathfrak m\) 分别化为 \(\bar g,\bar h\)。

\textbf{证明。}
这是一个逐次提升问题。证明思路与 \textit{Algebra III} 完成章节中的完备性工具一致，可视为 (III.Prop 8.5.1) 与完成 exactness 的标准应用。

先任取首一提升 \(g_1,h_1\in A[x]\)，使
\[
g_1\bmod \mathfrak m=\bar g,\qquad h_1\bmod \mathfrak m=\bar h.
\]
于是
\[
f-g_1h_1\in \mathfrak m A[x].
\]

现在归纳构造首一多项式 \(g_n,h_n\) 满足
\[
g_n\equiv g_{n-1}\pmod{\mathfrak m^{n-1}},\qquad
h_n\equiv h_{n-1}\pmod{\mathfrak m^{n-1}},
\]
并且
\[
f-g_nh_n\in \mathfrak m^nA[x].
\]

设 \(e_n=f-g_nh_n\in \mathfrak m^nA[x]\)。欲修正
\[
g_{n+1}=g_n+u,\qquad h_{n+1}=h_n+v,
\]
其中 \(u,v\in \mathfrak m^nA[x]\)，且次数限制使得首一性与次数不变。展开得
\[
f-g_{n+1}h_{n+1}
=e_n-g_nv-h_nu-uv.
\]
因为 \(uv\in \mathfrak m^{2n}A[x]\subseteq \mathfrak m^{n+1}A[x]\)（当 \(n\ge 1\)），所以模 \(\mathfrak m^{n+1}\) 只需解
\[
g_nv+h_nu\equiv e_n \pmod{\mathfrak m^{n+1}}.
\]
把它除以 \(\mathfrak m^n\) 并模 \(\mathfrak m\) 以后，变成在 \((A/\mathfrak m)[x]\) 中解
\[
\bar g\,\bar v+\bar h\,\bar u=\bar e_n.
\]

由于 \(\bar g,\bar h\) 互素，存在 Bézout 关系
\[
\alpha \bar g+\beta \bar h=1
\]
于 \((A/\mathfrak m)[x]\) 中。故对任意 \(\bar e_n\) 都可取
\[
\bar u=\beta \bar e_n,\qquad \bar v=\alpha \bar e_n
\]
解上式。把 \(\bar u,\bar v\) 提升到 \(u,v\in \mathfrak m^nA[x]\)，即可使误差提升到 \(\mathfrak m^{n+1}\)。

于是归纳可行，得到两个 Cauchy 列 \((g_n)\)、\((h_n)\)。由于 \(A\) 是 \(\mathfrak m\)-进完备的，系数逐项收敛，故存在极限
\[
g=\lim g_n,\qquad h=\lim h_n
\]
于 \(A[x]\) 中，并保持首一且
\[
g\bmod \mathfrak m=\bar g,\qquad h\bmod \mathfrak m=\bar h.
\]
又因
\[
f-g_nh_n\in \mathfrak m^nA[x]\qquad \forall n,
\]
令 \(n\to\infty\) 得
\[
f=gh.
\]
故 Hensel 引理成立。
\end{enumerate}

\section*{最后说明}
\phantomsection
\addcontentsline{toc}{section}{最后说明}

这一版相较于速记版，主要补了三类信息：
\begin{itemize}
  \item 省略步骤较多的题，改成可连续读下去的证明；
  \item 凡是明显调用了 \textit{Algebra III} 现成结论的地方，补了章节编号；
  \item 完成与局部化相关题目，把“为什么能局部化”“为什么核这样描述”尽量写明。
\end{itemize}

如果还要继续往下做，最自然的下一步是把其中几道偏难题再升级成“逐行证明版”，尤其是 5.15、10.3、10.9 这几题。

\end{document}
